% Chương 1. Người phụ trách chính: R1. Khoảng 1,5 đến 2 trang.
\chapter{Introduction}
\label{ch:intro}

\section{The oddball paradigm and the P300}
\fillpara{[R1, R3] Mô hình oddball: kích thích hiếm (target) xen giữa kích thích thường gặp (standard). Thành phần P300 (P3b): biên độ dương, vùng đỉnh trung tâm, khoảng 300 đến 600 ms; ý nghĩa chức năng (cập nhật ngữ cảnh, phân bổ chú ý). Phân biệt P3a (distractor) và P3b (target). Trích \citet{polich2007}.}

\section{Single-trial ERP classification}
\fillpara{[R1] Vì sao phân loại từng trial khó: tỉ số tín hiệu trên nhiễu thấp, khác biệt giữa người, mất cân bằng lớp (khoảng 1 target trên 6 standard). Ứng dụng: BCI, đánh giá chú ý. Vì sao cần balanced accuracy thay cho accuracy.}

\section{Classical machine learning and deep learning approaches}
\fillpara{[R4, R5] Classical: trung bình theo cửa sổ thời gian và shrinkage LDA; xDAWN cộng hình học Riemann \citep{rivet2009,barachant2012}. DL: EEGNet \citep{lawhern2018}, ShallowConvNet \citep{schirrmeister2017}. Câu hỏi còn mở: DL có hơn classical trên dữ liệu ERP nhỏ không.}

\section{Motor confound and explainability}
\fillpara{[R1] Sau khi chỉ giữ trial đúng, mọi target đều có bấm nút và không standard nào có, nên mô hình có thể học tín hiệu vận động thay cho P300. Vì sao cần giải thích mô hình (spatial pattern, saliency) và vì sao cần pattern thay cho weight \citep{haufe2014}.}

\section{Research questions}
\begin{itemize}
  \item \textbf{RQ1.} \fillin{dịch sang tiếng Anh: Kết quả phân loại single-trial tốt nhất cho frequent vs target?}
  \item \textbf{RQ2.} \fillin{dịch: Deep Learning có tốt hơn kỹ thuật ML tiêu chuẩn cho single-trial classification?}
  \item \textbf{RQ3.} \fillin{dịch: Đặc trưng mà mô hình DL khám phá có giải thích được không?}
\end{itemize}

\section{Report outline}
\fillpara{[R1] Một đoạn ngắn: Chương \ref{ch:data} mô tả dữ liệu, Chương \ref{ch:methods} phương pháp, Chương \ref{ch:results} kết quả, Chương \ref{ch:discussion} thảo luận, Chương \ref{ch:conclusion} kết luận.}
